% META: {"id": "1NSI-PM-EVAL-B-corrige", "chapitre": "1NSI-PROJET-METHODES", "type_objet": "corrige_evaluation", "evaluation_ref": "1NSI-PM-EVAL-B", "capacites": ["P-HIST-01", "BO-PREAMBULE-DEMARCHE-DE-PROJET", "BO-PREAMBULE-COMPETENCES-METHODE", "BO-PREAMBULE-COMPETENCES-ORALES"], "mode_creation": "ex_nihilo", "sources_inspiration": [], "status": "verified"}
% BEGIN-VERIFY
% evenements = [
%     (1936, "Concept de machine universelle", "Alan Turing"),
%     (1969, "Debuts du reseau Internet", "DARPA"),
%     (1991, "Premiere version publique de Python", "Guido van Rossum"),
% ]
% def evenements_entre(evenements, debut, fin):
%     return [e for e in evenements if debut <= e[0] <= fin]
% resultat = evenements_entre(evenements, 1960, 1991)
% assert [e[0] for e in resultat] == [1969, 1991]
% def acceder_element(liste, indice):
%     assert 0 <= indice < len(liste), "indice hors des bornes de la liste"
%     return liste[indice]
% assert acceder_element([1, 2, 3], 1) == 2
% def moyenne(notes, coefficients):
%     """Renvoie la moyenne ponderee des notes.
%
%     notes : liste de nombres, les notes obtenues.
%     coefficients : liste de nombres positifs, de meme longueur que notes.
%     Renvoie un nombre : la moyenne ponderee.
%     Precondition : la somme des coefficients est strictement positive.
%     """
%     assert sum(coefficients) > 0, "la somme des coefficients doit etre non nulle"
%     total = 0
%     poids = 0
%     for note, coefficient in zip(notes, coefficients):
%         total += note * coefficient
%         poids += coefficient
%     return total / poids
%
% assert moyenne([10, 20], [1, 1]) == 15
% assert moyenne([12, 8], [3, 1]) == 11
% for attendu in ("notes", "coefficients", "Renvoie", "Precondition"):
%     assert attendu in moyenne.__doc__
% try:
%     moyenne([10], [0])
% except AssertionError as e:
%     assert "coefficients" in str(e)
% else:
%     raise AssertionError("la precondition n'a pas ete declenchee")
% END-VERIFY
\begin{corrige}{1NSI-PM-EVAL-B-EX1}
\textbf{1.} Entre 1960 et 1991 : les débuts du réseau Internet (1969) et la première
version publique de Python (1991).

\textbf{2.} $\dfrac{3}{5}\times100=60{,}0\,\%$.

\textbf{3.} Avec seulement $2$ à $4$ élèves et un temps limité, un projet trop ambitieux
risque de ne \textbf{jamais aboutir} : l'équipe s'épuise sur des fonctionnalités
secondaires sans jamais disposer d'une version complète et testée. Un projet modeste mais
terminé et fonctionnel est pédagogiquement plus utile (et plus valorisant) qu'un projet
ambitieux resté à l'état d'ébauche.
\end{corrige}

\begin{corrige}{1NSI-PM-EVAL-B-EX2}
\textbf{1.} Le programme lève une \lstinline{IndexError}, avec le message
\lstinline{list index out of range}.

\textbf{2.}
\begin{python}
def acceder_element(liste, indice):
    assert 0 <= indice < len(liste), "indice hors des bornes de la liste"
    return liste[indice]
\end{python}

\textbf{3.} Une précondition explicite documente, directement dans le code, l'hypothèse
nécessaire (un indice valide) avec un message personnalisé et compréhensible. L'erreur
native \lstinline{IndexError} arrête aussi le programme, mais sans expliquer la
\textbf{cause métier} du problème (ici, un indice incohérent avec la taille de la liste)
aussi clairement qu'un message d'assertion rédigé pour ce cas précis.
\end{corrige}

\begin{corrige}{1NSI-PM-EVAL-B-EX3}
\textbf{1.} La docstring nomme le role, les paramètres avec leur type, la valeur renvoyee
et la precondition, et l'assertion la fait respecter :

\begin{python}
def moyenne(notes, coefficients):
    """Renvoie la moyenne ponderee des notes.

    notes : liste de nombres, les notes obtenues.
    coefficients : liste de nombres positifs, de meme longueur que notes.
    Renvoie un nombre : la moyenne ponderee.
    Precondition : la somme des coefficients est strictement positive.
    """
    assert sum(coefficients) > 0, "la somme des coefficients doit etre non nulle"
    total = 0
    poids = 0
    for note, coefficient in zip(notes, coefficients):
        total += note * coefficient
        poids += coefficient
    return total / poids
\end{python}

Sans cette precondition, un appel \lstinline{moyenne([10], [0])} echoue sur une division par
zero, dont le message ne dit rien du probleme réel.

\textbf{2.} Plan possible en trois points, chacun avec son argument principal :
\begin{itemize}
  \item \textbf{Le probleme et le besoin.} Argument : montrer a qui sert le programme et
  ce qu'il remplace, pour que le jury comprenne l'enjeu avant le code.
  \item \textbf{Une démonstration courte.} Argument : une execution réelle sur un jeu de
  donnees prepare prouve que le programme fonctionne, ce qu'aucune description ne remplace.
  \item \textbf{Un choix technique justifie.} Argument : par exemple, avoir choisi une
  \textbf{precondition} plutot qu'un \lstinline{try/except} parce qu'une somme de
  coefficients nulle revele une erreur de saisie en amont, qu'il vaut mieux signaler que
  rattraper silencieusement. C'est ce point qui \textbf{justifie} au lieu de decrire.
\end{itemize}
\end{corrige}
